% 06-discussion.tex — 讨论

\section{讨论}
\label{sec:discussion}

\subsection{Family Cap 的 ``均值降 / 峰值升'' 悖论}

小范围 sweep 的族数上限实验揭示了一个违反直觉的规律: maxFam=1 $\to$ 2 $\to$ 4 $\to$ 8 时, 均值从 0.6984 下降至 0.6306, 但全局峰值却由 maxFam=1 的 0.7999 升至 maxFam=8 的 0.8233。

这与神经架构搜索 (NAS) 的经验规律一致: 更大的搜索空间不一定带来更好的均值, 但峰值通常更高 \cite{liu2018darts, real2019regularized}。从激活族数分布看, maxFam=8 运行中 81.8\% 最终只保留 active$\le$3 的族, 且全场最高分 0.8233 本身仅 active=2 (即 6/8 族被关闭)。这说明 maxFam=k 的实际作用是把 ES 初始化为 k 组独立的 1 族起点, 而非真正同时激活 k 族进行联合优化。当某个种子恰好落在 ``对的多族组合'' 邻域时, 多族可突破单族天花板; 当所有种子都落在 ``远离任何好解'' 的邻域时, 多族反而稀释了最强信号。

\medskip\noindent\textbf{推论.} 在 Pareto 前沿上做 ``均值-峰值双目标优化'', 显式区分 ``均值稳健'' 与 ``峰值惊人'' 两类配置, 比折叠到单一标量 $maze\_quality$ 上更合理。

\subsection{两类失败模式的机制分离}

§5.3 的消融数据将失败模式清晰分为两类, 机制不同:

\noindent\textbf{(1) 预算受限型 (chebyshev-4, 均值 0.157).}
chebyshev-4 邻域严格包含 chebyshev-2, 理论上 chebyshev-2 的最优解是其搜索空间的子集, 但实际 chebyshev-4 得分低 0.59 分。原因不在 ``无解'', 而在 chebyshev-4 的 80 位 cells mask 占据 $2^{80}$ 种取值, 叠加 birth ($2^9$) + survive ($2^9$) + priority ($2^4$) 后, 单族总搜索空间 $= 2^{103} = 2^{80} \times 2^{23}$; $40 \times 60$ 的 $10^5$ 预算对其 $2^{103}$ 子空间过于稀疏, ES 在大量 ``无梯度平台'' 上原地踏步, 从未接近 chebyshev-2 风格的局部最优。这一类失败来自预算约束, 是高维搜索的固有挑战。

\noindent\textbf{(2) 表征受限型 (manhattan-1, 均值 0.421).}
manhattan-1 是 4 邻域 (von Neumann), 每个元胞每步最多感知 4 个邻居; 而走廊结构 (1-cell 宽, 长程连通) 需要 ``沿走廊方向存活 + 左右墙存活'' 的二阶条件, 4 邻域只能 ``凑巧'' 出现, 无法稳定涌现。与 manhattan-2 (8 邻域, 直接编码对角线存活条件) 差 0.358 分 --- 直接编码对角线存活条件的能力鸿沟。这一类失败来自模板本身的信息瓶颈, 增加预算无法解决。

\noindent\textbf{注意.}
chebyshev-4 和 manhattan-1 都得分低, 但分属不同类别。前者是搜索空间问题 (可增加预算或降低 mask 维度解决), 后者是表征能力问题 (必须换模板, 增加邻域半径或跳接结构才能解决)。将二者混为一谈会导致错误的调参方向: 在 manhattan-1 上加 10 倍预算不会让均值从 0.421 升到 0.75, 必须把模板换成 manhattan-2 才行。

\subsection{大范围 sweep 的含义: $40 \times 60$ 已接近饱和}

大范围 sweep 5/5 全部完成, 最高 0.8195 与小范围 sweep 最高 0.8233 几乎持平 --- 这一结果有重要意义。

\noindent\textbf{manhattan-2 / maxFam=8 在 $40 \times 60$ 网格上已接近该配置的理论上界。}将网格面积扩大 170 倍后, 峰值并未显著提升; 5 个随机种子中 s444 (0.8195) 与 s1111 (0.8138) 落在 ``好的吸引域'', s3333 (0.7909) 与 s2222 (0.7758) 落在 ``较差的吸引域'' (均低于 0.80), s5555 (0.8098) 接近均值。

5-run 均值 0.8020 与小范围 sweep 峰值 0.8233 差 0.021, 在搜索空间的方差范围内。top-2 均值 (s444 + s1111) 0.8166 与 0.8233 仅差 0.007, 说明 ``同配置在两个尺度上达到的水平'' 接近。

\noindent\textbf{为什么 $2^{1648}$ 的搜索空间在 $40 \times 60$ 上已被充分采样。}
manhattan-2 的 cells mask 是 80 位, 单族 $2^{103}$ 搜索空间; maxFam=8 时全染色体 $2^{824}$ 实际有效组合 (8 active 族 $\times$ 103 bits)。$10^5$ 预算相当于 $2^{17}$, 在 $2^{103}$ 上是 $2^{-86}$ 稀疏度; ES 用 $200 \times 500 = 10^5$ 步能稳定找到 $0.80$ 附近的局部最优, 但 5 个随机种子的标准差仍有 0.018, 表明 ``找到全局最优'' 在该预算下有 1/5 的概率。大范围 sweep 5 个种子中 s444 与 s1111 命中, s3333 与 s2222 未命中, 比例与 1/5 接近。

\noindent\textbf{推论.}
若要超过 0.8233, 需要:
\begin{itemize}[leftmargin=*, itemsep=0pt, topsep=0pt]
  \item 换 mask (manhattan-3 或 chebyshev-3 的均值与 manhattan-2 接近, 可能命中更好的吸引域);
  \item 增大预算 (pop $\times$ gens 从 $10^5$ 升到 $10^6$, 即大范围 sweep 的预算水平);
  \item 或两者都做。
\end{itemize}

本文未做新的 ``换 mask + 大范围'' sweep, 是后续工作的开放方向。

\subsection{局限与未来工作}

本文的局限有三个方面:
\begin{itemize}[leftmargin=*, itemsep=0pt, topsep=0pt]
  \item \textbf{权重选择}. $maze\_quality$ 的子度量权重 (如 $M_c = 0.40$, $M_t = 0.50$) 由经验决定, 理论上可通过元学习自动调优, 但本文未涉及。
  \item \textbf{15-pattern 完备性}. 15-pattern 数据集虽覆盖常见几何反例, 仍非完备; 可加入分形迷宫、Hilbert 曲线等拓扑反例。
  \item \textbf{浏览器预算}. ES 种群规模受 WebGPU dispatch 单次调用预算限制, $200 \times 500$ 已是浏览器端合理上限; 若部署到服务器端 GPU, 可尝试更大 pop $\times$ gens, 可能把 0.8233 继续上推。
\end{itemize}

FamilyMask 当前用 ``按优先级顺序遍历'' 执行 CA 步, 等价于 ``首个触发的族即决定下一态''; 这种 ``早退'' 语义在非 ordered ``多族投票'' 或 ``多族加权和'' 推广时需重新设计。
